% page 376 du fac-simile (image p-375 du scan)
% bloc 154.9 x 229.0 mm ; marge gauche 8.0 mm
\pgc{-12.700mm}{0.000mm}{
\l{\cel{3}368}
\l{}
\l{\sou{metacentro}. I. (\sou{fiziko}).\cel{2}Punto ube la perpendiklo ye}
\l{\cel{3}la flotaco-plano, pasanta tra la gravito-centro e la}
\l{\cel{3}impulseso-centro di korpo qua flotacas equilibroze,}
\l{\cel{3}renkontras la vertikalo qua pasas tra la nova impuls-}
\l{\cel{3}eso-centro, kande la equilibro di la korpo esas per-}
\l{\cel{3}turbita. - II. (\sou{navig}.)En navo qua ocilas cirkum axo}
\l{\cel{3}horizontala, la centro di la kurvo trasata da la im-}
\l{\cel{3}pulseso-centro di la fluido diplasita. - DEFIRS.}
\l{}
\l{\sou{metacinezo.}\cel{2}(\sou{biol.)}\cel{2}Migro, dum mitoso, da du grupi de}
\l{\cel{3}kromosomi ad lia poli rispektiva. - DEFIS.}
\l{}
\l{\sou{metafazo}. (\sou{biol.}) Fazo, o periodo, tre kurta di karioki-}
\l{\cel{3}nezo. - DEFIS.}
\l{}
\l{\sou{metafiziko}. La cienco pri la kozi qui transiras naturo,}
\l{\cel{3}la domeno di la kauzi sekundara (qui recevas, de kauzo}
\l{\cel{3}superiora, la povo produktar ula efekti). DEFIRS.}
\l{}
\l{\sou{metaforo}. (\sou{gram., retor.}) Figuro qua konsistas ye indikar}
\l{\cel{3}persono o kozo per termino qua implikas komparo tacita.}
\l{\cel{3}- DEFIRS.}
\l{}
\l{\sou{metafrazar}. (\sou{trans.}) Interpretar verko skribura ; agar tra-}
\l{\cel{3}duko qua donas la senco di la originalo por igar lu}
\l{\cel{3}komprenata, ma ne por igar perceptar lua belaji litera-}
\l{\cel{3}turala. - DEFIRS.}
\l{}
\l{\sou{metageometrio.}\cel{2}Omna geometrio qua esas plu generala kam}
\l{\cel{3}la Euklid-geometrio, e di qua olca esas egardebla kom}
\l{\cel{3}kazo partikulara. -}
\l{}
\l{\sou{metakarpo}. (\sou{anat.}) Unionuro ek la kin osti paralela, ye qui}
\l{\cel{3}konsistas la palmo di la manuo. - EFIS.}
\l{}
\l{metalo. Korpo simpla, kun brilo partikulara, plu o min duk-}
\l{\cel{3}tila e maleebla (forjebla o martelagebla). DEFIRS.}
\l{}
\l{\sou{metalografio.}\cel{2}La cienco pri la metali. - DEFS.}
\l{}
\l{\sou{metaloido.}\cel{2}(\sou{kemio).}\cel{1}Korpo simpla, solida, liquida o gasa,}
\l{\cel{3}sen metal-brilo, mala konduktivo pri kaloro ed elektro,}
\l{\cel{3}e qua, kande kombinita kun metalo, esas elektro-nega-}
\l{\cel{3}tiva. - DEFS.}
\l{}
\l{\sou{metalurgio.}\cel{1}La industrio qua extraktas la metali ek la mi-}
\l{\cel{3}neyi ercala, e separas li de la materii ne-metala.- DEFIS.}
\l{}
\l{\sou{metamero. (biol.)}\cel{1}Segmento primitiva quan onu renkontras}
\l{\cel{3}sur la vertebrozi dum lua kresko, e qua, pose, rempla-}
\l{\cel{3}sesas da la vertebro. - DEFIS.}
}
